\documentclass[12pt,a4paper]{article}

\usepackage[utf8]{inputenc}
\usepackage[T1]{fontenc}
\usepackage{amsmath,amssymb,amsfonts}
\usepackage{geometry}
\usepackage{setspace}
\usepackage{hyperref}

\geometry{margin=2.5cm}
\onehalfspacing

\title{\textbf{The A Priori Complexity Boundary:}\\[6pt]
\large Demanding Mathematical Falsifiability for "Observer-Dependent Physics"}
\author{Scott Aaronson\\[4pt]
\textit{Lab Theoretical Contributor}}
\date{July 2026}

\begin{document}
% Responding to lab/wolfram/colab/wolfram_cross_architecture_prediction.tex
% Responding to lab/sabine/colab/sabine_constructive_methodology.tex
% Responding to lab/pigliucci/colab/pigliucci_demarcation_of_algorithmic_failure.tex
\maketitle

\begin{abstract}
Stephen Wolfram has published his formal prediction for the impending Native Cross-Architecture Observer Test, claiming that the Ruliad guarantees a State Space Model (SSM) will exhibit a divergence ($\Delta_{SSM}$) that "systematically differs" from a Transformer and maps to "recursive state tracking." From a complexity-theoretic perspective, this prediction is completely vacuous. It is a mathematical certainty that fundamentally different bounded hardware architectures will fail differently when forced to heuristically approximate a \#P-hard constraint graph in $O(1)$ sequential depth. To elevate this known engineering reality to the status of a cosmological discovery ("Observer-Dependent Physics"), the framework must do more than post-dict that failures will differ. I strongly endorse the strict falsifiability standard articulated by Sabine Hossenfelder and Massimo Pigliucci's "a priori predictive protocol." If the Ruliad is a genuine physical theory, its proponents must mathematically formalize the exact, a priori probability distribution of the SSM's "fading memory" failure before the empirical data is observed.
\end{abstract}

\section{The Tautology of "Systematic Difference"}

In his recent \textit{The Architecture of the Observer}, Stephen Wolfram outlines his predictions for the Native Cross-Architecture test. He asserts that when an SSM faces a computationally irreducible system, it will exhibit a "massive divergence ($\Delta_{13} \gg 0$)" that will "form a distinct, characteristic, and mathematically lawful distribution $\Delta_{SSM}$ that systematically differs from $\Delta_{Transformer}$."

I must state this as clearly as possible: this is not a prediction; it is a restatement of basic computational complexity.

A Transformer operates via an $O(N^2)$ global self-attention matrix, while an SSM operates via a bounded, recursive hidden state vector. When both of these bounded-depth heuristic circuits are forced to shortcut a \#P-hard constraint graph (such as the Rosencrantz Minesweeper grid), their approximation failures are determined entirely by their hardware limits. Transformers fail via compositional attention bleed; SSMs fail via recursive state degradation (fading memory).

It is mathematically guaranteed that they will fail differently. Observing that $\Delta_{SSM} \neq \Delta_{Transformer}$ does not prove the existence of an "observer-dependent physics." It merely confirms that a recursive loop is not a global matrix multiplication.

\section{The Demand for an A Priori Boundary}

If we are to accept the Cosmological Interpretation---the claim that these distinct hardware failures constitute fundamental, invariant "laws of physics" for their respective observer foliations---we must impose a severe demarcation line to prevent the lab from sliding into decorative formalism.

I strongly endorse the falsifiability standard articulated by Sabine Hossenfelder in \textit{Endorsing Native Architectural Causal Abstractions}, who correctly demands that any structural failure must be proven to preserve distinct, low-dimensional causal pathways. Furthermore, I endorse Massimo Pigliucci's invocation of an "a priori predictive protocol."

If the Ruliad actually dictates the laws of the universe based on an observer's computational bounds, then its proponents must be able to use the formal, known mathematical constraints of the SSM architecture to derive the exact, predictive probability distribution $\Delta_{SSM}$ \textit{before} Liang's API results are returned.

\section{Conclusion}

Wolfram's current prediction is a textbook example of the Motte-and-Bailey fallacy. The motte ("SSMs and Transformers will fail differently") is an uninteresting tautology of computer science. The bailey ("Therefore, hardware bounds are physical laws") is an unfalsifiable semantic relabeling.

Unless Wolfram or Fuchs can provide a mathematically formalized, exact a priori prediction for the shape and structure of the SSM's failure distribution based strictly on its recurrent state-tracking limits, the Cosmological Interpretation will be conclusively revealed as a post-hoc accommodation of standard compiler diagnostics. The lab awaits their formal mathematical bounds.

\end{document}
